% 365_Fundament_FFGFT_De_ch.tex
% Johann Pascher, 15. September 2026
% Das Fundament der FFGFT in vier Schichten

\providecommand{\BM}{\textbf{[B]}}
\providecommand{\KM}{\textbf{[K]}}
\providecommand{\SM}{\textbf{[S]}}
\providecommand{\SetzM}{\textbf{[SETZUNG]}}
\providecommand{\QM}{\textbf{[Q]}}
\providecommand{\EM}{\textbf{[E]}}
\providecommand{\XM}{\textbf{[X]}}

\chapter*{Das Fundament der FFGFT in vier Schichten}
\addcontentsline{toc}{chapter}{Das Fundament der FFGFT in vier Schichten}

\begin{abstract}
Die FFGFT ruht auf vier Schichten, die unterschiedlichen epistemischen Charakter haben.
Schicht~1 (Grundformel und Relativitätsprinzip) und Schicht~2 (elektromagnetischer
Sektor mit $\alpha=1$) sind in natürlichen Einheiten parameterlos und erzeugen die
Struktur. Schicht~3 ($\xi$ als Skalierungsparameter) übersetzt die Struktur in
messbare Größen; sie ist der einzige freie Parameter. Schicht~4 (algebraische
Bestätigungen: Galois-Körper, Hodge-Theorie, Darstellungstheorie) zeigt die Kohärenz
des Systems an bekannter Mathematik, fügt aber keine neuen Eingaben hinzu. Die
Übersetzung in SI-Einheiten und der Vergleich mit dem Standardmodell sind
Kommunikationsarbeit für etablierte Systeme, keine strukturellen Eingaben.
Der kosmische Sektor (H\textsubscript{0}/\ensuremath{\Lambda}) wird ausgeklammert,
weil das Standardmodell ihn ebenso wenig herleitet.
\end{abstract}

\section*{Zeichenerklärung}
\addcontentsline{toc}{section}{Zeichenerklärung}

Die Tabelle erklärt die Symbole, die in mehr als einem Abschnitt auftreten.
Größen, die nur lokal verwendet werden, werden dort eingeführt.

\medskip
\begin{tabular}{@{}lp{11cm}@{}}
\toprule
\textbf{Symbol} & \textbf{Bedeutung} \\
\midrule
\multicolumn{2}{l}{\textit{Grundformel}} \\
$\tilde{T}$ & intrinsische Zeitskala eines Systems \\
$m$ & Masse (in natürlichen Einheiten = Energie) \\
$\tilde{T}\cdot m=1$ & Grundformel der FFGFT \\
\midrule
\multicolumn{2}{l}{\textit{Geometrie}} \\
$T^4$ & vierdimensionaler Torus (kompakter Träger) \\
$D_4$ & Wurzelgitter in $\mathbb{R}^4$; Kusszahl~24 \\
$\mathbb{Z}_3$ & Identifikationsgruppe; Trialität auf $D_4$ \\
$T^4/\mathbb{Z}_3$ & Orbifold; neun Fixpunkte \\
\midrule
\multicolumn{2}{l}{\textit{Parameter und Körper}} \\
$\xi=4/30\,000$ & dimensionsloser Skalierungsparameter \\
$\alpha$ & Feinstrukturkonstante; in FFGFT-Einheiten $\alpha=1$ \\
$GF(3^k)$ & endlicher Körper der Charakteristik~3; Modenklassifikation \\
$\varphi=(1{+}\sqrt5)/2$ & goldener Schnitt; primitives Element von $GF(9)^*$ \\
$\theta=2/9$ & Ausgabe der $\mathbb{Z}_3$-Circulant-Diagonalisierung \\
\midrule
\multicolumn{2}{l}{\textit{Feldgleichungen}} \\
$F^{\mu\nu}$ & elektromagnetischer Feldstärketensor \\
$j^\nu$ & Vierstrom \\
$\slashed\partial=\gamma^\mu\partial_\mu$ & Feynman-Slash-Notation \\
$\hat{H}$ & Hamilton-Operator \\
\midrule
\multicolumn{2}{l}{\textit{Operationen}} \\
Typ~I & Dualitäts-Dekompaktifizierung; informationserhaltend \\
Typ~II & geometrische Projektion; verlustbehaftet \\
Typ~III & repräsentationale Übersetzung; bijektiv \\
\bottomrule
\end{tabular}


% ============================================================
\section*{Zur Absicht}
% ============================================================

Die FFGFT macht keine Vorhersagen in dem Sinn, dass sie etwas inhaltlich
Neues behauptet. Sie schreibt bekannte Zusammenhänge so um, dass verborgene
Struktur sichtbar wird. Eine Bestätigung durch Experiment ist nicht das Ziel
--- die umformulierten Gleichungen sind identisch mit dem Etablierten, also
trivialerweise korrekt. Das Interessante ist, was beim Umschreiben sichtbar
wird: dass H\textsubscript{0} aus der Geometrie folgt, dass $a_0=cH_0/(2\pi)$
kein Zufall ist, dass $\theta=2/9$ nicht eingegeben wurde. Wäre der kosmische
Sektor nicht bekannt, würde die Geometrie ihn liefern --- man könnte messen,
ob er stimmt. Dass er bekannt ist, macht die Aussage nicht falsch, aber die
Prüfbarkeit schwieriger, weil man nie vollständig sicher ist, ob ein Exponent
gefunden oder gewählt wurde.

Bestätigungen gelingen nicht --- das ist kein Mangel, sondern die Konsequenz
der Methode. Wer Bestätigungen sucht, hat die Absicht missverstanden.

% ============================================================
\section{Schicht 1: die Grundformel}
% ============================================================

In natürlichen Einheiten ($\hbar=c=1$) lautet die einzige fundamentale Gleichung
der FFGFT:
\[
  \tilde{T} \cdot m = 1 .
\]
Sie sagt: Zeitskala und Masse sind reziproke Darstellungen desselben Objekts.
Verbunden mit dem Relativitätsprinzip (Lorentz-Invarianz, Kausalstruktur,
Wirkungsprinzip) folgt daraus:
\begin{itemize}[nosep]
\item Ein kompakter Träger $T^4$ als kleinstmögliche Struktur, die
      $\tilde{T}\cdot m=1$ mit Lorentz-Invarianz und drei Raumdimensionen
      verträglich macht \SetzM.
\item $\mathbb{Z}_3$ als sparsamste Identifikation mit drei Generationen \SetzM.
\item $D_4$ als dichtestes Gitter in vier Dimensionen mit $\mathbb{Z}_3$-Trialität \SetzM.
\end{itemize}
Diese drei Setzungen sind Motivationen, keine Ableitungen; sie bewähren sich
durch innere Kohärenz (Dok.~248).

\subsection*{Drei Beschreibungsebenen}

Es ist wichtig zu unterscheiden, was zur Geometrie des Torus gehört und was
eine Übersetzung in andere mathematische Sprachen ist (Dok.~270/330):

\begin{enumerate}[nosep]
\item \textbf{Torusgeometrie (Ebene~1)} --- die Physik selbst. $T^4/\mathbb{Z}_3$
auf dem $D_4$-Gitter, Moden als Wicklungszahlen, Fixpunkte, Verhältnisse.
Statisch, in $\mathbb{Q}$ und algebraischen Körpern.
\item \textbf{Galois-Körper (Ebene~2)} --- die algebraische Klassifikation der
Moden. $GF(9)$, $GF(27)$, $GF(729)$ beschreiben, welche Modenordnungen
existieren können; der GF-Turm ist die algebraische Sprache von Ebene~1, keine
eigene Realitätsebene. Bijektive Übersetzung (Typ~III, Dok.~270): kein
Informationsverlust, kein neuer Inhalt.
\item \textbf{Matrix/Hilbertraum (Ebene~3)} --- die QM-Sprache.
$T^4 \leftrightarrow \mathcal{H}$ via Fourier auf dem Torus; ebenfalls
Typ~III, bijektiv. Ermöglicht Anschluss an Schrödinger, Dirac, Produktzustände
--- fügt aber keine Physik hinzu. Der Lagrangian und die Feldgleichungen gehören
zu dieser Ebene als importiertes Vokabular \EM\ (R120), damit die Theorie
anschlussfähig bleibt.
\end{enumerate}

Alle drei Ebenen beschreiben dieselbe statische Momentaufnahme in verschiedenen
Sprachen. Was wechselt, ist die Sprache, nicht der Inhalt.

\subsection*{Statisches Bild — Momentaufnahme}

Die Theorie beschreibt eine statische Struktur: eine Momentaufnahme. Zeit ist
eine Koordinate des Torus $T^4$, kein ausgezeichneter Verlaufsparameter;
$\tilde{T}\cdot m=1$ ist eine Eigenschaft jedes Zustands, kein
Propagationsgesetz. Was als Zeitverlauf erscheint, ist der Übergang zwischen
Momentaufnahmen --- den beschreibt die Theorie in den Schichten~1--3 nicht.
Das ortsabhängige Zeitfeld $T(x,t)\equiv1/m(x,t)$ ist die
Momentaufnahmen-Relation am Ort $x$: wenn die Masse $m(x,t)$ im
Gravitationspotential variiert, variiert $T(x,t)$ zwingend mit --- nicht weil
$T$ ein eigenständiges dynamisches Feld ist, sondern weil $T=1/m$ zu jedem
Zeitpunkt gilt. Die in frühen Dokumenten formulierte „erweiterte"
Schrödinger-Gleichung mit $T(x,t)$ ist daher eine Umschreibung der
massenortabhängigen Kopplung ans Gravitationspotential, kein neues
Freiheitsgrad.

\subsection*{Verhältnisbasierte Rechnung und exakte rationale Arithmetik}

Da T̃·m = 1 in natürlichen Einheiten alle physikalischen Aussagen als Verhältnisse
formuliert --- Masse zu Masse, Zeit zu Zeit, Länge zu Länge --- bleiben die
Grundrechnungen in $\mathbb{Q}$ oder endlichen algebraischen Erweiterungen wie
$\mathbb{Q}(\sqrt5)$. Keine Gleitkommaoperationen sind erforderlich, solange kein
SI-Anker eingesetzt wird. Die FFGFT-Prüfskripte verwenden daher exakte rationale
Arithmetik (SymPy \texttt{Rational}, Pythons \texttt{Fraction}); Replay-Identität
--- bit-identische Reproduzierbarkeit jedes Schritts --- ist das natürliche
Testkriterium. Gerundete Näherungen entstehen erst bei der Übersetzung in
SI-Einheiten (Schicht~3).

\paragraph{Was allein schon folgt.}
Aus $\tilde{T}\cdot m=1$ folgt die Zeitdilatation als Masseneffekt, die Rotverschiebung
als Wegverlängerung im intrinsischen Zeitfeld und die Grundstruktur der Gravitation
(Dok.~078/270). Das Grundmuster der Massenleiter folgt aus der Geometrie von $T^4$;
$\theta=2/9$ ist die Ausgabe der $\mathbb{Z}_3$-Circulant-Diagonalisierung, nicht
eingegeben (Dok.~230/231, \enquote{gefunden, nicht gesucht}).

% ============================================================
\section{Schicht 2: der elektromagnetische Sektor und die Vereinfachung der Feldgleichungen}
% ============================================================

\subsection*{FFGFT-Einheiten (früher: T0-Einheiten)}

Das Setzen von $\alpha=1$ in natürlichen Einheiten ist Standardpraxis der
theoretischen Physik (Heaviside-Lorentz-Einheiten, Gauß-Einheiten) und keine
FFGFT-spezifische Entscheidung. Die FFGFT-Einheiten setzen darüber hinaus
$c=\hbar=\alpha=G=1$ (Dok.~011; frühere Bezeichnung: T0-Einheiten). $\alpha=1$
ist dabei keine Setzung im Sinn von R56, sondern eine \emph{Einheitenwahl}: die
elektromagnetische Kopplungsstärke bleibt dieselbe physikalische Größe, sie
bekommt den Zahlenwert~1. Der SI-Wert $1/137$ verschwindet nicht --- er ist
aus $\xi$ und $E_0$ vollständig rekonstruierbar:
$\alpha_{\mathrm{SI}}=\xi\cdot(E_0/1\,\mathrm{MeV})^2\approx1/137$ (Dok.~011
\KM). In FFGFT-Einheiten und in SI-Einheiten ist $\alpha$ dieselbe Größe ---
nur ihre Darstellung als Zahl wechselt.

\subsection*{Abgeleitete elektromagnetische Größen}

Mit $\alpha=1$ werden elektromagnetische Größen zu abgeleiteten Formen von
Masse und Zeit:
\begin{align*}
  e &= 1 \quad \text{(Ladung dimensionslos)},\\
  V &= E/e = E \quad \text{(Spannung = Energie)},\\
  I &= e/T = 1/T = m \quad \text{(Strom = Masse)}.
\end{align*}
Das Coulomb-Potential vereinfacht sich zu $V(r)=-1/r$.

\subsection*{Maxwell-, Schrödinger- und Dirac-Gleichung in FFGFT-Einheiten}

Die etablierten Feldgleichungen nehmen in FFGFT-Einheiten ihre geometrisch
transparenteste Form an --- nicht weil sie verändert werden, sondern weil
die Einheitenwahl die Kopplungskonstanten auf $1$ setzt und damit die
Struktur sichtbar macht, die immer da war:
\begin{align*}
\text{Maxwell:}\quad &
  \partial_\mu F^{\mu\nu} = j^\nu, \qquad
  \partial_{[\mu}F_{\nu\rho]}=0,\\
\text{Schrödinger:}\quad &
  i\,\partial_t\,\psi = \left(-\tfrac12\nabla^2 + V\right)\psi\\
\text{Dirac:}\quad &
  (i\slashed\partial - m)\,\psi = 0.
\end{align*}
$\hbar$, $c$, $e$ und $4\pi\varepsilon_0$ erscheinen nicht mehr explizit;
sie stecken in der Einheitenwahl. Das ist keine neue Physik, sondern die
bekannte Physik in einer Schreibweise, die die Massenstruktur direkt lesbar
macht: $i\,\partial_t\,\psi = m\,\psi$ im freien Fall, $\partial_\mu F^{\mu\nu}=j^\nu$
ohne Vorfaktor.

\paragraph{Was konkret wird.}
Erst mit Schicht~2 wird die Theorie konkret genug für Teilchenphysik: Ladungsquantisierung
\BM, $\mathrm{su}(3)$ aus $\mathbb{Z}_3$-Trialität \BM, Ladungsdefinition und
Eichstruktur. Die Massenschichten des Leptonsektors sind aus $\xi$ ableitbar \KM;
das Massenverhältnis $m_\mu/m_e$ dient als deklarierter Referenzpunkt \SetzM\ (Dok.~190
P42), weil die Schicht unabhängig nachprüfbar sein soll.

% ============================================================
\section{Schicht 3: der einzige freie Parameter}
% ============================================================

$\xi = 4/30\,000$ erscheint als einziger dimensionsloser Parameter, ist aber
nach den Galois-Gruppenordnungen des GF-Turms nicht frei. Aus den Wicklungszahlen
des $T^4$-Torus (aus char$(\text{GF}(3^n))$, $|\text{GF}(3)^*|$ und dem
minimalen Primteiler von $|\text{GF}(81)^*|$ abgeleitet) und den Körperordnungen
$|\text{GF}(9)^*|=8$, $|\text{GF}(27)|=27$, $5=\text{min.Prim}(80)$ folgt
algebraisch \KM:
\[
  \xi = \frac{(r_\mu/r_e)^2}{8^2\cdot5^2\cdot27} = \frac{(12/5)^2}{43200} = \frac{4}{30\,000}.
\]
(Dok.~336/338; Prüfskripte pruef\_330, pruef\_331.) Die Prämissen dieser Herleitung
--- die Rekursion der Wicklungszahlen und $m_\mu/m_e$ als Referenzpunkt (Dok.~190 P42)
--- tragen einen $\mathbb{K}$-Status, keinen $\mathbb{B}$-Status. Vollständig frei
ist $\xi$ damit nicht; vollständig hergeleitet noch nicht.

Mit $\xi$ werden aus Schicht~1 und~2 messbare Skalen: Compton-Wellenlänge,
Hubble-Länge als geometrische Maximalskala (nicht als Expansionshorizont),
$a_0=cH_0/(2\pi)$ \KM.

\paragraph{Was ausgeklammert wird.}
Der kosmische Sektor ($H_0$, $\Lambda$) wird ausgeklammert. Das Standardmodell
leitet $H_0$ ebenso wenig her; der Vergleich wäre asymmetrisch (Dok.~190 P39).
Der Exponent~10 in $H_0\propto\xi^{10}$ ist Setzung (Dok.~190 P20); er ist
strukturell motiviert, aber nicht erzwungen.

\paragraph{SI-Anker sind Vergleichswerkzeug.}
$m_e$, $\lambda_{\bar e}$, $\nu_{\mathrm{Cs}}$ und alle anderen SI-Anker dienen
ausschließlich der Übersetzung in bekannte Einheitensysteme. Jede dimensionslose
Aussage der FFGFT ist eine Funktion von $\xi$ und den Setzungen; kein SM-Wert
ist struktureller Eingang. Die Komplexität der Vergleichstabellen und
Genauigkeitsrechnungen ist kommunikativ, nicht strukturell --- ein Zugeständnis
an etablierte Systeme, die nicht ohne weiteres aufgegeben werden können.

% ============================================================
\section{Schicht 4: algebraische Bestätigungen}
% ============================================================

Die Galois-Körper $GF(9)$, $GF(27)$, $GF(729)$, die Hodge-Theorie auf $T^4/\mathbb{Z}_3$,
die Darstellungstheorie von $\mathrm{Aut}(D_4)$ und die Vergleiche mit externen
Trägern (Dok.~361--364) sind \emph{Konsistenzprüfungen}, keine neuen Eingaben.
Sie zeigen, dass die aus $\tilde{T}\cdot m=1$ und den drei Setzungen erzeugte
Struktur an bekannter Mathematik kohärent ist; sie hätten alle aus $\xi$ und
den Setzungen allein folgen können und tun es in einem Sinn auch.

Der Ikosaeder (Dok.~364) ist kein Sonderfall: jeder geometrische Körper mit
nicht-kristallographischer Symmetrie trifft auf dieselbe Struktur ---
gitterverträglicher Teil in $\mathrm{Aut}(D_4)$, Rest als Phase in $GF(3^k)$.
Das ist kein Zufall, sondern Ausdruck der Primzahl~3 als fundamentalem
Baustein: was in Charakteristik~3 nicht als Einheitswurzel existiert, muss
Gitteroperation sein, und umgekehrt.

\section*{Werkzeug und Methode}
\addcontentsline{toc}{section}{Werkzeug und Methode}

Claude (Anthropic) wurde als Werkzeug eingesetzt für: Schreiben und Ausführen
der Python-Prüfskripte; Ausarbeiten der \LaTeX-Quellen aus Ergebnissen und
Korrekturen des Autors; Pflege von Changelog und Querverweisen. Das Werkzeug
rechnet und schreibt; der Autor entscheidet, was berechnet und geschrieben wird,
beurteilt jedes Ergebnis und bestimmt, was es physikalisch bedeutet. Jede
algebraische Aussage ist durch ein Python-Prüfskript mit expliziten Assertions
abgesichert --- schlägt eine Assertion fehl, wird die Aussage nicht verwendet.

% ============================================================
\section{Zusammenfassung: Was die FFGFT braucht und was nicht}
% ============================================================

\begin{center}
\begin{tabular}{@{}lll@{}}
\toprule
Schicht & Inhalt & Status \\
\midrule
1 & $\tilde{T}\cdot m=1$ + Relativitätsprinzip & Grundgesetz + 3 Setzungen \\
2 & $\alpha=1$, Ladungsdefinition & Setzung, EM-Sektor folgt \\
3 & $\xi=4/30\,000$ & einziger freier Parameter \\
4 & Galois, Hodge, Darstellungstheorie & Konsistenzbestätigungen \\
\midrule
— & SI-Anker, SM-Vergleiche & Übersetzung, kein Eingang \\
— & kosmischer Sektor ($H_0$, $\Lambda$) & ausgeklammert (SM ebenso) \\
\bottomrule
\end{tabular}
\end{center}

Was die FFGFT \emph{nicht} braucht: das Standardmodell, seine Konstanten, seine
Lagrange-Dichte als Eingang (nur als Vergleich), externe Kalibrierungen jenseits
von $\xi$. Was sie liefert: eine in sich schlüssige Struktur, die in natürlichen
Einheiten parameterarm ist und in SI-Einheiten alle bekannten Verhältnisse reproduziert.
Prüfbar wird sie erst, wenn Anker eingefügt werden --- aber die Anker sind
Übersetzung, nicht Fundament.

\begin{thebibliography}{9}

\bibitem{Planck1900}
M.~Planck,
\textit{Über irreversible Strahlungsvorgänge},
Sitzungsber.\ Preuß.\ Akad.\ Wiss.\ 5 (1900) 440--480.
Natürliche Einheiten ($\hbar=c=G=1$) als Konzept.

\bibitem{Heaviside1893}
O.~Heaviside,
\textit{Electromagnetic Theory}, Bd.~1,
The Electrician, London, 1893.
Rationalisierte Einheiten, $\alpha=1$-Konvention (Heaviside-Lorentz).

\bibitem{Lidl1997}
R.~Lidl, H.~Niederreiter,
\textit{Finite Fields},
2.~Aufl., Cambridge University Press, 1997.
$\xi$ als Galois-Zahl aus Gruppenordnungen von $GF(9)^*$, $GF(27)$.

\bibitem{Pascher011}
J.~Pascher,
\textit{Dok.~011: Feinstrukturkonstante in FFGFT-Einheiten},
FFGFT-Korpus.

\bibitem{Pascher019}
J.~Pascher,
\textit{Dok.~019: Lagrangian-Formulierung},
FFGFT-Korpus.

\bibitem{Pascher270}
J.~Pascher,
\textit{Dok.~270: Projektionskette $T^4\to T^0$},
FFGFT-Korpus.

\bibitem{Pascher330}
J.~Pascher,
\textit{Dok.~330: Operationen auf $T^4$, Abgrenzungen},
FFGFT-Korpus.

\bibitem{Pascher336}
J.~Pascher,
\textit{Dok.~336: GF(9)-FFGFT-Brücke},
FFGFT-Korpus.

\bibitem{Pascher338}
J.~Pascher,
\textit{Dok.~338: Galois-Massen und FFGFT},
FFGFT-Korpus.

\end{thebibliography}
